\documentclass[12pt,a4paper]{ctexart}
\usepackage[a4paper,margin=2.7cm]{geometry}
\usepackage{amsmath,amssymb}
\usepackage{booktabs}
\usepackage{array}
\usepackage{enumitem}
\usepackage{setspace}
\usepackage{titlesec}
\usepackage{hyperref}
\usepackage{fontspec}
\setmainfont{Times New Roman}

\setstretch{1}
\titleformat{\section}{\centering\bfseries\large}{}{0pt}{}
\titleformat{\subsection}{\bfseries\normalsize}{}{0pt}{}

\begin{document}

\begin{center}
    {\LARGE \textbf{Outline of Modern Chinese History}}
\end{center}

\section{I. Old Democratic Revolution}

\begin{enumerate}[leftmargin=2em]
    \item \textbf{1840--1842: The First Opium War}: It marked the beginning of modern Chinese history.
    \item \textbf{1841: \textit{Illustrated Gazetteer of the Maritime Countries}}: Lin Zexu was known as one of the first Chinese to ``open his eyes to the world.''
    \item \textbf{1843: \textit{Illustrated Treatise on the Maritime Kingdoms}}: Wei Yuan proposed ``learning advanced techniques from foreigners in order to control them.''
    \item \textbf{1851--1864: The Taiping Heavenly Kingdom Movement}: Its representative documents included \textit{The Land System of the Heavenly Dynasty} and \textit{A New Essay on Aid in Administration}.
    \item \textbf{1856--1860: The Second Opium War}: It further deepened China's semi-colonial and semi-feudal condition.
    \item \textbf{1861: The Self-Strengthening Movement}: It advocated self-strengthening and the pursuit of wealth.
    \item \textbf{1894--1895: The Sino-Japanese War of 1894--1895}: It ended with the Treaty of Shimonoseki.
    \item \textbf{1898: The Reform Movement of 1898}: It was a bourgeois political reform movement.
    \item \textbf{1898: \textit{On Evolution}}: Yan Fu introduced ideas such as ``natural selection'' and ``survival of the fittest.''
    \item \textbf{1899--1900: The Boxer Movement}: Its slogan was ``Support the Qing and destroy the foreign powers.''
    \item \textbf{1900--1901: The Invasion by the Eight-Power Allied Forces}: It culminated in the Boxer Protocol.
    \item \textbf{1905: Founding of the Revolutionary Alliance}: The Tongmenghui was founded in Tokyo, Japan.
    \item \textbf{1911: The Railway Protection Movement}: It was especially prominent in Sichuan.
    \item \textbf{1911: The Revolution of 1911}: It overthrew the Qing dynasty and brought down the feudal monarchy.
    \item \textbf{1912: Establishment of the Nanjing Provisional Government}: The \textit{Provisional Constitution of the Republic of China} was the first constitution of a bourgeois republic in China.
    \item \textbf{1912: Abdication of the Qing Emperor}: It marked the collapse of the feudal autocratic monarchy.
    \item \textbf{1913: The Second Revolution}: It was triggered by the assassination of Song Jiaoren.
    \item \textbf{1915: The New Culture Movement}: Chen Duxiu launched \textit{Youth Magazine}.
    \item \textbf{1915: The National Protection Movement}: It opposed Yuan Shikai's attempt to restore monarchy.
    \item \textbf{1917: The Constitutional Protection Movement}: It opposed the Beiyang warlords' refusal to restore the provisional constitution and parliament.
    \item \textbf{1917: The October Revolution in Russia}: It brought Marxism to China.
\end{enumerate}

\section{II. New Democratic Revolution}

\begin{enumerate}[leftmargin=2em]
    \item \textbf{1919: The May Fourth Movement}: Its slogans included democracy and science.
    \item \textbf{1921: The First National Congress of the Communist Party of China}: It marked the founding of the Communist Party of China.
    \item \textbf{1922: The Second National Congress of the Communist Party of China}: It put forward an anti-imperialist and anti-feudal program of national revolution.
    \item \textbf{1923: The Third National Congress of the Communist Party of China}: It decided on cooperation between the Communist Party and the Kuomintang.
\end{enumerate}

\section{III. The Great Revolution}

\begin{enumerate}[leftmargin=2em]
    \item \textbf{1924: The First National Congress of the Kuomintang}: It marked the first cooperation between the Kuomintang and the Communist Party and the formation of the New Three People's Principles.
    \item \textbf{1925: The Fourth National Congress of the Communist Party of China}: It put forward the worker-peasant alliance.
    \item \textbf{1926: The Northern Expedition}: Its slogan was ``Down with imperialism and the warlords.''
    \item \textbf{1927: The Fifth National Congress of the Communist Party of China}: It proposed striving for proletarian leadership, establishing revolutionary national power, and carrying out land revolution.
    \item \textbf{1927: The April 12 Counterrevolutionary Coup}: It marked the collapse of the first cooperation between the Kuomintang and the Communist Party and the failure of the Great Revolution.
\end{enumerate}

\section{IV. Agrarian Revolutionary War}

\begin{enumerate}[leftmargin=2em]
    \item \textbf{1927: The July 15 Counterrevolutionary Coup}: It was associated with Wang Jingwei.
    \item \textbf{1927: The Nanchang Uprising}: It fired the first shot of armed resistance against Kuomintang reactionary rule.
    \item \textbf{1927: The August 7 Meeting}: It emphasized that ``political power grows out of the barrel of a gun.''
    \item \textbf{1927: The Autumn Harvest Uprising}: It was led by Mao Zedong.
    \item \textbf{1927: The Ning-Han Split Reconciled}: It involved Nanjing and Wuhan, associated with Chiang Kai-shek and Wang Jingwei.
    \item \textbf{1927: The Sanwan Reorganization}: It established the Party's absolute leadership over the army.
    \item \textbf{1928: The Jinggangshan Meeting of Forces}: It opened the correct revolutionary road of encircling the cities from the countryside and seizing power by armed force.
    \item \textbf{1928: The Northeastern Flag Replacement}: Zhang Xueliang declared obedience to the National Government.
    \item \textbf{1928: The Sixth National Congress of the Communist Party of China}: It was held in Moscow.
    \item \textbf{1928: Mao Zedong's Two Important Articles}: These included \textit{Why Can China's Red Political Power Exist?} and \textit{The Struggle in the Jinggang Mountains}.
    \item \textbf{1929: The Gutian Meeting}: It established the principles of Party building ideologically and army building politically.
    \item \textbf{1930: Mao Zedong's Another Two Important Articles}: These included \textit{A Single Spark Can Start a Prairie Fire} and \textit{Oppose Book Worship}.
    \item \textbf{1931: The September 18 Incident}: It marked the beginning of Japanese imperialist aggression against China.
    \item \textbf{1934: The Beginning of the Long March}: It followed the failure of the fifth campaign against encirclement and suppression.
    \item \textbf{1935: The Zunyi Meeting}: It was a turning point in the history of the Communist Party of China.
    \item \textbf{1935: The Wayaobao Meeting}: It proposed the anti-Japanese national united front.
    \item \textbf{1936: The End of the Long March}: The main forces met in Huining, Gansu.
    \item \textbf{1936: The Xi'an Incident}: It was associated with Zhang Xueliang and Yang Hucheng.
\end{enumerate}

\section{V. The War of Resistance Against Japanese Aggression}

\begin{enumerate}[leftmargin=2em]
    \item \textbf{1937: The July 7 Incident}: It marked the beginning of full-scale Japanese invasion.
    \item \textbf{1937: Statement Recognizing the Legal Status of the Communist Party}: It marked the beginning of the second cooperation between the Kuomintang and the Communist Party.
    \item \textbf{1937: The Battle of Shanghai}: It was the first large-scale battle between China and Japan.
    \item \textbf{1937: The Battle of Taiyuan}: It featured the victory at Pingxingguan by the Communist-led forces.
    \item \textbf{1938: The Battle of Xuzhou}: It included the victory at Taierzhuang.
    \item \textbf{1938: \textit{On Protracted War}}: Mao Zedong emphasized that the army and the people were the foundation of victory.
    \item \textbf{1939: \textit{The Communist}}: It summarized the three magic weapons: the united front, armed struggle, and Party building.
    \item \textbf{1941: The Yan'an Rectification Movement}: It included important texts such as \textit{Reform Our Study}, \textit{Rectify the Party's Style of Work}, and \textit{Oppose Stereotyped Party Writing}.
    \item \textbf{1941: The Southern Anhui Incident}: The Kuomintang slandered the Communist Party.
    \item \textbf{1945: The Seventh National Congress of the Communist Party of China}: It established Mao Zedong Thought as the Party's guiding ideology.
    \item \textbf{1945: Japan's Unconditional Surrender}: It marked the victory of the Chinese People's War of Resistance Against Japanese Aggression.
\end{enumerate}

\section{VI. The War of Liberation}

\begin{enumerate}[leftmargin=2em]
    \item \textbf{1945: The Double Tenth Agreement}: It was signed in Chongqing.
    \item \textbf{1945: Chiang Kai-shek tore up the Double Tenth Agreement}: It marked the breakdown of the second cooperation between the Kuomintang and the Communist Party.
    \item \textbf{1947: The Battle of Menglianggu}: It crushed the Kuomintang's plan for a decisive battle in central Shandong.
    \item \textbf{1947: The Army's Advance into the Dabie Mountains}: It marked the beginning of the strategic counteroffensive.
    \item \textbf{1948: The Three Major Campaigns}: These were the Liaoshen, Huaihai, and Pingjin campaigns.
    \item \textbf{1949: The Second Plenary Session of the Seventh Central Committee of the Communist Party of China}: It was held in Xibaipo.
    \item \textbf{1949: The Crossing-the-Yangtze Campaign}: Its slogan was ``Cross the Yangtze River and liberate all of China.''
    \item \textbf{1949: Founding of the People's Republic of China}: It marked the great victory of the New Democratic Revolution.
\end{enumerate}

\end{document}